\documentclass[11pt]{article}
\usepackage{amsmath,amssymb}
\title{Abstract Mathematics: Introduction to Metric Spaces}
\author{Synthetic playground test prepared by Sol}
\date{}
\begin{document}
\maketitle

\section*{One concept: a metric space}
Let $X$ be a nonempty set. A metric on $X$ is a function
\[
 d:X\times X\longrightarrow [0,\infty)
\]
such that, for every $x,y,z\in X$:
\begin{enumerate}
 \item $d(x,y)=0$ if and only if $x=y$;
 \item $d(x,y)=d(y,x)$;
 \item $d(x,z)\leq d(x,y)+d(y,z)$.
\end{enumerate}
Nonnegativity is included in the codomain: $d(x,y)\geq 0$.
The pair $(X,d)$ is called a metric space. Its points need not be
points in the Euclidean plane; the set and its metric specify the space.

\section*{A small illustration of the same concept}
For $X=\mathbb{R}$, the usual metric is $d(x,y)=|x-y|$.
For $x=1$, $y=3$ and $z=4$,
\[
 d(1,3)=2,\qquad d(3,4)=1,\qquad d(1,4)=3\leq 2+1.
\]
This calculation illustrates the triangle inequality for one triple;
it does not prove the metric axioms for all real numbers.
The general definition above, including its universal quantifiers,
is the concept; this example does not replace it.

\section*{Scope and provenance}
This is original synthetic test material, not an extract from Gyile's
course notes. A drawing of creatures or connections may serve as a
mnemonic for the definition. Their positions, edges or colours do not
by themselves specify a metric, verify its axioms or prove a theorem.
\end{document}
